\subsection{General cyclic CM types and the sum of sharp directional slopes}

\begin{theorem}[DP40: cyclic displays and the alternating maximum]\label{DP40}
Let $h\ge2$, $q=p^h$, $s=1/(p-1)$, and let $A$ be a nonempty proper subset of $\Z/h\Z$, of cardinality $m$. Over $R=W(k)$ take a Witt module with cyclic basis $e_i$, normal decomposition
\[
 T=\bigoplus_{i\in A}W(R)e_i,\qquad
 L=\bigoplus_{i\notin A}W(R)e_i,\qquad Q=L\oplus IT,
\]
and arbitrary units $u_i\in W(R)^\times$. Define
\[
 Fe_i=u_ip^{\epsilon_i}e_{i+1},\qquad
 \epsilon_i=1-1_A(i),\qquad
 V^{-1}e_i=u_ie_{i+1}\quad(i\notin A),
\]
with the defining display rule on $IT$. This gives a formal $p$-divisible group $H_{A,u}$ of height $h$ and dimension $m$. Its actual Cartier dual is the complementary display $H_{A^c,u^{-1}}$.

For $j\in A$, put
\[
 r_j=\frac{\sum_{i\in A}p^{[j-i]_h}}{q-1},\qquad
 c_j=s-r_j=\frac{\sum_{i\notin A}p^{[j-i]_h}}{q-1},
\]
where $[a]_h$ is the least nonnegative residue. In the Teichm\"uller coordinates $X_j$, the normalized integral invariant directions satisfy, coefficientwise,
\begin{equation}\label{eq:cyclicmixed}
 \vp\left(\prod_{j\in A}D_j^{n_j}f\right)
 \ge\max\left\{0,\sum_{j\in A}c_jn_j-
                   \frac{\sum_{j\in A}s_p(n_j)}{p-1}\right\}.
\end{equation}
Each slope is sharp, even with the arbitrary unit arrows:
\[
 \vp(D_j^{q^r}X_j(0))=c_j(q^r-1)\qquad(r\ge0).
\]
For unit arrows equal to one the exact logarithm is
\[
 L_A=\sum_{n\ge0}\frac{M^{\circ n}}{p^n},\qquad
 M_j=X_{\operatorname{prev}(j)}^{p^{g_j}},
\]
where $g_j$ is the positive cyclic gap from the preceding selected index to $j$.

Across the group and its actual dual define
\[
 \gamma_i(A)=\frac{\sum_{j:\,1_A(j)\ne1_A(i)}p^{[i-j]_h}}{q-1},
 \qquad N_k=|A\mathbin\triangle(A+k)|,\qquad
 S_h(A)=\sum_i\gamma_i(A).
\]
Then
\begin{equation}\label{eq:cyclicmax}
 S_h(A)=\frac{\sum_{k=1}^{h-1}N_kp^k}{q-1}
 \le h c_2,\qquad
 \sum_{i\in A}\gamma_i=\sum_{i\notin A}\gamma_i=\frac{S_h(A)}2.
\end{equation}
Equality in the first inequality holds exactly when $h$ is even and $A$ is alternating. More precisely,
\[
 S_h(A)\le \frac{4m(h-m)}h\,c_2.
\]
For nonconstant types equality in this last inequality again requires an even alternating type. These are equally weighted sums, not componentwise bounds or global theta costs.
\end{theorem}

\begin{proof}
The linearized normal map $V^{-1}\oplus F:L\oplus T\to P$ is a monomial matrix of units, hence invertible. Complete cycles accumulate $h-m>0$ powers of $p$ under $F$, and the complementary cycle accumulates $m>0$. The corresponding nilpotence conditions hold over each $R/p^n$ and are compatible under reduction. Zink's normal-decomposition construction, p.~129, and Theorem~5, p.~131, give a formal $p$-divisible group of height $\operatorname{rank}P=h$; its tangent module $P/Q$ has rank $m$. On the dual basis the evaluation pairing forces
\[
 F'e_i^*=u_i^{-1}p^{1-\epsilon_i}e_{i+1}^*,\qquad
 T'=L^*,\quad L'=T^*.
\]
Indeed $\langle Fe_i,F'e_j^*\rangle=p\delta_{ij}$, and the other identities follow on the normal generators. Both displays are connected. Zink's Definition~19/(22), Proposition~121 and Theorem~122 identify this complementary display as the actual Cartier dual \cite{Zink}.

The section $\sum_{i\in A}[X_i]e_i$ has identity tangent map to $P/Q$, so gives integral formal coordinates. Write $z_n=w_n(z)/p^n$ for divided Witt ghosts. With unit arrows one, let $B$ denote the Frobenius matrix. The resolution operator on these ghosts is $B\operatorname{shift}-1$, and its logarithm is
\[
 \operatorname{pr}_T\sum_{n\ge0}B^nz_n.
\]
It vanishes on the relation module by telescoping and has identity differential. For arbitrary units, let $B_n$ have nonzero entries $w_n(u_i)p^{\epsilon_i}$. Each $w_n(u_i)$ is a unit, and the operator and logarithm become
\[
 z_n\longmapsto B_nz_{n+1}-z_n,\qquad
 \operatorname{pr}_T\left(z_0+\sum_{n\ge1}B_0\cdots B_{n-1}z_n\right).
\]
Here is an explicit completion justification. On a degree truncation with nilpotent ideal $N^D=0$, put $C=\floor{\log_p(D-1)}$ if $D>1$. Integral Witt vectors satisfy
\[
 z_n=\sum_{r=0}^{\min(n,C)}
       \frac{a_{n-r}^{p^r}}{p^r};
\]
thus their divided ghosts have uniformly bounded denominators. Every complete matrix cycle contributes $h-m$ powers of $p$ times units. The sums therefore converge in the completed integral resolution, and their telescoping boundaries tend to zero. The particular Teichm\"uller sections are finite on every degree truncation. This is a rational calculation with the already constructed integral quotient, not an application of a display theorem over a field with $p$ inverted.

From one selected index to the next, a path of length $g_j$ has exactly one Frobenius arrow without a factor $p$. Its Teichm\"uller ghost denominator is $p^{g_j}$, leaving precisely a unit divided by $p$. A path through $a$ selected gaps has coefficient $\eta/p^a$, with $\eta\in R^\times$. Each starting index and ordinary degree specifies a unique path, so there is no cancellation. With all units one, these paths give exactly the asserted formula for $L_A$.

The weights satisfy
\begin{equation}\label{eq:cyclicedge}
 p^{g_j}r_{\operatorname{prev}(j)}-r_j=1.
\end{equation}
Indeed multiplication of the defining numerator by $p^{g_j}$ advances each cyclic distance to $j$; only the term belonging to $j$ wraps, changing $1$ to $q$. Adjoin $\pi^{q-1}=p$ and scale coordinate $j$ by $\alpha_j=\pi^{(q-1)r_j}$. A path of one selected gap then has valuation zero. A path of $a\ge2$ selected gaps has scaled valuation
\[
 1+p^{g_1}+p^{g_1+g_2}+\cdots+
       p^{g_1+\cdots+g_{a-1}}-a>0,
\]
with gaps ordered backwards from the output. Thus the contracted logarithm is integral, has identity differential, and reduces to
\[
 \overline{L'}=\id+M_\eta,\qquad
 (M_\eta)_j=\bar\eta_jX_{\operatorname{prev}(j)}^{p^{g_j}},
 \quad\bar\eta_j\ne0.
\]
Formal inversion is integral. Since $M_\eta$ is additive in characteristic $p$ and raises ordinary degree, its inverse reduction is
\[
 \overline{E'}=\sum_{a\ge0}(-1)^aM_\eta^{\circ a}.
\]
In component $j$, with only input coordinate $j$ retained, degree $q^r$ comes from exactly $mr$ arrows and has a nonzero coefficient. Hence
\[
 \vp([t^{q^r}]E_j(te_j))=-(q^r-1)r_j.
\]
For any integral $f$, contraction bounds the coefficient of $\mathbf t^{\mathbf n}$ in $f(E(\mathbf t))$ below by $-\sum r_jn_j$. Multiplication by $\prod n_j!$ and Legendre's formula give \eqref{eq:cyclicmixed} at zero. Apply it to the integral translated function to obtain the coefficientwise statement. For $f=X_j$, $n=q^r$, the exact coefficient gives
\[
 \vp(D_j^nX_j(0))=\vp(n!)-(n-1)r_j=(n-1)c_j.
\]
This proves sharpness and works also at $p=2$.

It remains to prove the sum comparison. At displacement $k$ there are $N_k$ pairs of opposite type. The number with output in $A$ is
$|A|-|A\cap(A+k)|$; the same formula counts those with output in $A^c$. Each is $N_k/2$, proving both sum identities. Also $N_k=N_{h-k}$.

Set $x_i=1$ on $A$ and $x_i=-1$ otherwise. For the circulant operator
\[
 (Wx)_i=(q-1)^{-1}\sum_{k=1}^{h-1}p^kx_{i-k}
\]
we have
\[
 S_h(A)=\frac12(hc_h-x^TWx).
\]
At the discrete frequencies $\theta_\ell=2\pi\ell/h$, the symmetric part of $W$ has eigenvalues
\[
 \lambda_\ell=\frac{\sum_{k=1}^{h-1}p^k\cos(k\theta_\ell)}{q-1}
 =\frac{p\cos\theta_\ell-1}{p^2-2p\cos\theta_\ell+1}
       -\frac1{q-1}.
\]
The rational expression follows from the finite geometric sum using
$e^{ih\theta_\ell}=1$; it is not asserted for arbitrary real frequencies.
As a function of $t=\cos\theta_\ell$ its derivative is
$p(p^2-1)/(p^2-2pt+1)^2>0$. Consequently
\[
 \lambda_\ell\ge\lambda_*=-\frac1{p+1}-\frac1{q-1},
\]
with equality only at $\theta_\ell=\pi$. For normalized Fourier coefficients
$\widehat x_\ell=h^{-1/2}\sum_jx_je^{-2\pi i\ell j/h}$, Parseval gives
$\sum|\widehat x_\ell|^2=h$. Hence
\[
 S_h(A)\le\frac h2(c_h-\lambda_*)=hc_2.
\]
Equality forces Fourier support only at $h/2$, so $h$ is even and
$x_j=\pm(-1)^j$. Conversely those two words give equality. Explicitly,
\[
 hc_2-S_h(A)=\frac{p(p-1)}{2(p+1)}
 \sum_{\ell=0}^{h-1}
 \frac{1+\cos\theta_\ell}{p^2-2p\cos\theta_\ell+1}
 |\widehat x_\ell|^2.
\]
Finally retain the zeroth eigenvalue $\lambda_0=c_h$ and
$|\widehat x_0|^2=(2m-h)^2/h$ instead of bounding it by $\lambda_*$. This gives
$S_h(A)\le hc_2-c_2(2m-h)^2/h=4m(h-m)c_2/h$.
Equality permits only modes $0,h/2$; a nonconstant $\{\pm1\}$ word of that form is alternating. Constant types, if included, have sum zero and equality in this cardinality bound.
\end{proof}

The display hypotheses have a useful exact CM interface. Suppose an integral action of an unramified degree-$h$ order splits the display into free rank-one eigenspaces, Frobenius cycles those eigenspaces, and the Hodge filtration has $Q_i=IP_i$ on $A$ and $Q_i=P_i$ off $A$. Choosing generators, commutation with the action forces the arrow positions above, and invertibility of the normal map forces their coefficients to be units. Thus DP40 applies. This argument does not infer those integral hypotheses from height alone.

For $h=6$ and $A=\{0,1,2\}$, direct counting gives
\[
 (N_1,\ldots,N_5)=(2,4,6,4,2),\qquad
 (\gamma_0,\gamma_1,\gamma_2)
 =(\gamma_3,\gamma_4,\gamma_5)
 =\frac{(p,p^2,p^3)}{(p-1)(p^3+1)}.
\]
Therefore
\[
 S_6(A)=6c_2-\frac{4p(p-1)}{p^3+1},\qquad
 \frac{S_6(A)}{6c_2}=\frac{p^2+p+1}{3(p^2-p+1)}<1.
\]
Nevertheless $\gamma_2-c_2=p/(p^3+1)>0$, which explicitly refutes a componentwise interpretation. DP22's CM type $\{1,2,3\}$ for $y^2=x^7-1$ gives this word on the Frobenius orbits
$(1,3,2,6,4,5)$ when $p\equiv3\bmod7$ and
$(2,3,1,5,4,6)$ when $p\equiv5\bmod7$.
The finite check of all binary words through $h=10$ at $p=2,3,5,11$ verifies 8,184 cases of the sum identities, inequalities and equality classifications; it supplements the proof above.

